\scanpage{117}\mspage{62}
à examiner, qui n'ont pas l'air du tout de provenir des foncteurs du diagramme p.~8 ou de quelque composé (ça ne donne que des équivalences canoniques ou bien les foncteurs $i,j$ de page 51\unsure{)|,}, quand on va des cat. internes \unsure{au|en} diagramme vers les catégories \unsure{bordantes}). Le plus joli est sans doute de les regarder à partir de $\mathrm{Cat}_{/S}$, puisque c'est là seulement qu'on les a tous les deux à la fois
\diagram{\begin{tikzcd}[ampersand replacement=\&,column sep=small,row sep=small]
 \mathrm{Cat}_{/S}(W_S^d)^{-1}\approx\mathrm{HOT}_W(S)\arrow[dr,"\unsure{\rho|p}"] \& {}\\
 {} \& \mathrm{Cat}_{/S}W_S^{-1}\approx\mathrm{HOT}_W^{\unsure{\mathrm{lc}|\ell\mathrm{c}}}(S)\\
 \mathrm{Cat}_{/S}(W_S^g)^{-1}\approx\mathrm{HOT}_W^\circ(S)\arrow[ur,"\sigma"] \& {}
\end{tikzcd}}

Par contre, les foncteurs (\unsure{qu'on a envie|que j'ai envie} d'appeler ``foncteurs d'\unsure{induction|inclusion}'') $i,j$ de la page 51, en sens inverse de $\rho,\sigma$, sont ceux qui proviennent, par localisation $W_S^d$ ou $W_S^g$ selon les cas, d'un des foncteurs du diagramme page 8 (ou d'un composé de tels foncteurs), allant d'une des dix catégories internes (où toutes les $W$-équivalences envisagées coïncident) vers l'une des sept catégories \unsure{bordantes}, sauf quand cette dernière est $\mathrm{Cat}_{/S}$, et détermine déjà avec quelle $W$-équivalence, $W_S^d$ ou $W_S^g$, on va travailler. La détermi%
% archive sheet number: 116
